\chapter{Laços}

\section{\texttt{for}}

Itera sobre uma sequência. A variável de iteração é visível dentro do bloco e
\textit{persiste} após o laço com o último valor assumido:

\begin{lstlisting}[caption={for sobre range}]
for i in 1..4 {
  display i    // 1, 2, 3  (4 não incluído)
}
display i    // 3 — variável persiste após o laço
\end{lstlisting}

\subsection*{Range exclusivo e inclusivo}

\hayashi{} segue a mesma convenção de Rust:

\begin{lstlisting}
for i in 1..3  { }    // 1, 2      (exclusivo: não inclui o fim)
for i in 1..=3 { }    // 1, 2, 3   (inclusivo: inclui o fim)
\end{lstlisting}

Ranges também são expressões que avaliam para listas — ver Seção~\ref{sec:ranges}.

\subsection*{Iterando sobre listas}

\begin{lstlisting}
let frutas = ["maçã", "banana", "uva"]

for fruta in frutas {
  display fruta
}
\end{lstlisting}

\subsection*{Descartando a variável de iteração}

Quando o índice não importa:

\begin{lstlisting}
for _ in 1..3 {
  display "repetição"
}
\end{lstlisting}

\section{\texttt{while}}

Executa enquanto a condição for verdadeira:

\begin{lstlisting}[caption={while básico}]
let n = 1

while n <= 5 {
  display n
  n += 1
}
\end{lstlisting}

\begin{aviso}
  \hayashi{} não detecta laços infinitos. Um \lstinline|while true { }| sem
  \lstinline{break} executará indefinidamente. Use \texttt{Ctrl-C} para interromper.
\end{aviso}

\section{\texttt{break} e \texttt{continue}}

\lstinline{break} interrompe o laço imediatamente. \lstinline{continue} avança para
a próxima iteração:

\begin{lstlisting}[caption={break e continue}]
for i in 1..10 {
  if i == 3 { continue }    // pula o 3
  if i == 7 { break }       // para no 7
  display i
}
// exibe: 1, 2, 4, 5, 6
\end{lstlisting}

\section{Laços com DataFrames}

\lstinline{for} pode iterar sobre listas produzidas por funções de dados:

\begin{lstlisting}[caption={Iterando sobre grupos}]
load "vendas.csv" as df

for uf in ["RS", "SP", "RJ"] {
  let sub = filter(df, estado == uf)
  display f"{uf}: {count(sub)} obs"
}
\end{lstlisting}

\section{\texttt{parallel for}}

Variante concorrente do \texttt{for}. As iterações executam em paralelo
entre threads; o valor de retorno de cada iteração (última expressão ou
\texttt{return} explícito) é coletado em uma lista.

\textbf{Como expressão} (preferido --- o resultado vai para uma variável
nomeada):

\begin{lstlisting}[caption={parallel for como expressão}]
let resultados = parallel for t in tickers {
  let sub = filter(df, ticker == t)
  nrow(sub)
}
// resultados é uma lista de valores nrow
\end{lstlisting}

\textbf{Como statement} (o resultado é armazenado na variável de
iteração):

\begin{lstlisting}[caption={parallel for como statement}]
parallel for t in tickers {
  let sub = filter(df, ticker == t)
  nrow(sub)
}
// t agora contém uma lista de valores nrow
\end{lstlisting}

O parâmetro opcional \texttt{threads=N} limita o número de threads
trabalhadoras (padrão: todas as CPUs disponíveis):

\begin{lstlisting}[caption={Limite de threads}]
let resultados = parallel for t in tickers, threads=4 {
  load "data.db" as sub, query=f"SELECT * FROM p WHERE t='{t}'"
  nrow(sub)
}
\end{lstlisting}

\subsection*{Modelo de isolamento}

Cada thread recebe seu próprio interpretador com um \textbf{snapshot} do
ambiente externo (apenas valores send-safe são capturados). Isso tem
consequências importantes:

\begin{itemize}
  \item \textbf{Leituras do escopo externo funcionam}: variáveis definidas
    antes do \texttt{parallel for} são visíveis dentro do corpo.
  \item \textbf{Escritas no escopo externo são silenciosamente descartadas}:
    reatribuir uma variável externa dentro do corpo não tem efeito sobre
    o escopo externo. A mutação existe apenas no ambiente local da thread
    e é perdida quando a thread termina.
\end{itemize}

\begin{lstlisting}[caption={Mutações são descartadas}]
let n = 0
parallel for i in 1..3 {
  n = 10          // sem efeito sobre n externo
  i * 2
}
print(n)          // ainda 0
\end{lstlisting}

\begin{aviso}
  Trate o corpo como uma \textbf{função pura}: leia variáveis externas,
  compute um resultado, e faça \texttt{return} dele. Não confie em
  efeitos colaterais para propagar valores para fora.
\end{aviso}

\subsection*{O que pode e não pode ser capturado}

\begin{center}
\begin{tabular}{lll}
\toprule
\textbf{Tipo} & \textbf{Capturável?} & \textbf{Motivo} \\
\midrule
\texttt{Int}, \texttt{Float}, \texttt{Str}, \texttt{Bool} & Sim & Primitivos, copiados por valor \\
\texttt{List}, \texttt{Dict} (send-safe) & Sim & \texttt{Arc}, leitura compartilhada \\
\texttt{DataFrame}, \texttt{Series} & Sim & \texttt{Arc}, operações são funcionais \\
\texttt{UserFn} & Sim & \texttt{Arc} \\
Resultados de modelos (\texttt{ols}, \texttt{iv}, \texttt{rolling}, \ldots) & \textbf{Não} & Contêm \texttt{Rc}, não são \texttt{Send} \\
\bottomrule
\end{tabular}
\end{center}

Resultados de modelos do escopo externo não podem ser capturados. Compute
modelos \textbf{dentro} do corpo:

\begin{lstlisting}[caption={Modelos devem ser computados dentro}]
// ERRADO: m contém Rc, será rejeitado
let m = ols(y ~ x, df)
parallel for i in 1..10 { predict(m, newdata) }

// CORRETO: computar dentro do corpo
parallel for i in 1..10 {
  let sub = filter(df, grupo == i)
  let m = ols(y ~ x, sub)
  tidy(m)
}
\end{lstlisting}

\subsection*{Pular iterações e agregação}

Use \texttt{return nil} para pular uma iteração. Entradas \texttt{nil} são
mantidas na lista de resultados (em ordem). Combine com \texttt{rbind()}
para agregar DataFrames por iteração --- \texttt{rbind} ignora
silenciosamente entradas \texttt{nil}:

\begin{lstlisting}[caption={parallel for + rbind}]
let resultados = parallel for i in 0..n, threads=8 {
  let t = tickers[i]
  if t == "SPY" { return nil }
  let df_t = calcular_betas(t)
  df_t
}
let todos = rbind(resultados)
\end{lstlisting}
